\author{Brett Reynolds}
\title{ভাষার ভূদৃশ্য: Language Landscapes}
\subtitle{ইংরেজি শিক্ষকদের জন্য ইংরেজি গঠনের নির্দেশিকা [বাংলা সংস্করণ]}
% \ISBNdigital{}
% \ISBNhardcover{}
% \BookDOI{}
\typesetter{Brett Reynolds}
% \proofreader{}
% \lsCoverTitleSizes{51.5pt}{17pt}
\renewcommand{\lsSeries}{tbls}
%\newcommand{\lsSeriesTitle}{Textbooks in Language Sciences}
%\newcommand{\lsSeriesColor}{lsYellow}
\renewcommand{\lsSeriesNumber}{16}
\BackBody{ভাষা মুখস্থ করার তালিকা নয়; ভাষা অনুসন্ধান করার এক ভূদৃশ্য।

ভালো মাঠ-প্রকৃতিবিদ শুধু নাম দেন না। কীভাবে জিনিসগুলি আচরণ করে, কোথায় একত্র হয়, কেন বদলে যায়, তিনি তা পর্যবেক্ষণ করেন। এই বই ইংরেজিকে একইভাবে পড়তে শেখায়। নিয়মপুস্তক ইংরেজি সম্পর্কে কী দাবি করে, তার চেয়ে আসল ভাষাভাষীরা কী করেন, সেটাই এখানে মুখ্য। তাত্ত্বিক কাঠামো \textit{The Cambridge Grammar of the English Language}-এর উপর দাঁড়ালেও উপস্থাপনাটি সহজ করা হয়েছে। কিন্তু মান কমানো হয়নি।

অধ্যায়গুলি ধ্বনি ও শব্দ থেকে গুচ্ছ, উপবাক্য, এবং বক্তব্যের দিকে এগোয়। পরিষ্কার ইংরেজি উদাহরণের সঙ্গে অন্য ভাষার ছোট ছোট জানালা খোলা হয়। প্রতিটি অধ্যায়ে একই কাজ করা হয়: ছক দেখা, প্রমাণ দিয়ে পরীক্ষা করা, কোথায় ছক ভাঙে তা দেখা, এবং সেই ভাঙন আমাদের কী শেখায় তা জিজ্ঞাসা করা। শেষে আপনি এমন ব্যবহারযোগ্য ব্যাকরণ-জ্ঞান পাবেন, যার সাহায্যে পাঠ্যবই মূল্যায়ন করা, শিক্ষার্থীর বাক্য কেন একটু অস্বাভাবিক লাগছে তা বোঝানো, এবং প্রশ্ন উঠলে নিজের ব্যাখ্যা রক্ষা করা সম্ভব।

আপনি শিক্ষক-প্রশিক্ষণের পথে থাকুন বা ইতিমধ্যে শ্রেণিকক্ষে পড়ান, লাভ একই। অযৌক্তিক নিয়মের উপর ভর করা বন্ধ করে ভাষা-ব্যবস্থাকে নিজে দেখতে শেখা।}
